\section{Product Strategy and Commercial Model}
\label{sec:business}

\noindent\emph{Note on figures.} The quantities in this section are modelling
assumptions with stated bases, not measured market research. They are included
because a design document without a commercial frame cannot be evaluated for
proportionality --- an engineering programme costing \$4.4\,M must be justified
against a plausible return. Every assumption is labelled and each should be
replaced with measured data before it is used for a funding decision.

\subsection{Positioning}

\EM{} occupies a position that is currently empty: professional-grade imaging
science delivered through a consumer mobile interface. The two adjacent
positions are both crowded.

\begin{itemize}
\item \textbf{Below}: consumer film-look camera applications. Large audience,
low price, minimal technical differentiation, high churn, competition on
aesthetics and marketing rather than capability.
\item \textbf{Above}: desktop film emulation plugins. Small audience, high
price, deep technical differentiation, low churn, competition on fidelity.
\end{itemize}

The wager is that a meaningful population sits between them: people who care
enough about image quality to notice the difference and who do not have, or do
not want to use, a desktop workflow for phone captures. This population is the
P1 and P3 personas of Section~\ref{sec:requirements}.

The defensibility argument is that the gap has persisted not because it is
uninteresting but because closing it requires a combination of skills ---
sensitometry, GPU engineering, and consumer product design --- that rarely
co-occurs, and because the work does not decompose into a weekend project. That
same difficulty is the moat.

\subsection{Differentiation That Survives Copying}

Features can be copied; the following cannot easily be:

\begin{enumerate}
\item \textbf{The fitted profile library.} A stock profile is the output of
acquiring physical film, exposing calibrated wedges, processing, measuring, and
fitting. That pipeline is a capital and time investment, and it compounds: each
additional stock is cheaper than the last because the tooling amortizes.
\item \textbf{The validation apparatus.} The test suite of
Section~\ref{sec:validation} is what allows the model to be changed safely. A
competitor who copies the equations without the validation cannot iterate on
them.
\item \textbf{Architectural performance.} The $1.9\times$ from LUT baking and
tile fusion (Table~\ref{tab:optlog}) is architectural. A competitor building on
a filter-chain framework cannot reach it without rewriting.
\item \textbf{Vocabulary credibility.} If the printer lights behave like printer
lights, professionals notice and say so. That credibility is earned slowly and
is difficult to fake, because the people who would be fooled are not the people
whose opinion moves the market.
\end{enumerate}

\subsection{Commercial Model}

A hybrid model is proposed:

\begin{table}[htbp]
\caption{Proposed Pricing Structure}
\label{tab:pricing}
\begin{center}
\renewcommand{\arraystretch}{1.15}
\begin{tabularx}{\columnwidth}{@{}l l >{\raggedright\arraybackslash}X@{}}
\toprule
\textbf{Tier} & \textbf{Price} & \textbf{Contents} \\
\midrule
Free & --- & Capture, develop, export at reduced resolution. Three negative stocks, one print stock. Full pipeline, not a crippled one. \\
Pro & \$5.99/mo or \$39.99/yr & Full stock library, full resolution export, recipes, batch, sidecar interchange, all future stocks. \\
Lifetime & \$129.99 & Pro, permanently. Offered for the first twelve months only. \\
\bottomrule
\end{tabularx}
\end{center}
\end{table}

The free tier deliberately includes the \emph{complete} pipeline rather than a
reduced one. The product's argument is that physical simulation is
perceptibly better; a free tier that removes halation or grain removes the
argument. What the free tier limits is breadth (stocks) and output (resolution),
which are the dimensions a serious user needs and a casual user does not.

The lifetime option exists because a meaningful fraction of the P1/P2 population
actively resents subscriptions for creative tools, and because early lifetime
revenue funds the expensive Phase 4 profile work. Its twelve-month window
prevents it from becoming a permanent drag on recurring revenue.

\subsection{Unit Economics Model}

\begin{table}[htbp]
\caption{Illustrative Year-2 Model}
\label{tab:unitecon}
\begin{center}
\renewcommand{\arraystretch}{1.15}
\begin{tabularx}{\columnwidth}{@{}>{\raggedright\arraybackslash}X r l@{}}
\toprule
\textbf{Quantity} & \textbf{Value} & \textbf{Basis} \\
\midrule
Downloads, year 2 & 850{,}000 & Assumption \\
Free-to-paid conversion & 3.2\% & Assumption; category range 1--6\% \\
Paying users acquired & 27{,}200 & Derived \\
Annual-plan share & 62\% & Assumption \\
Blended first-year ARPU & \$34.10 & Derived \\
Gross revenue, year 2 & \$927{,}000 & Derived \\
App Store commission & 15\% & Small business programme \\
Net revenue, year 2 & \$788{,}000 & Derived \\
Year-2 cost & \$2{,}498{,}000 & Table~\ref{tab:cost} \\
\midrule
Year-2 net & $-$\$1{,}710{,}000 & Derived \\
\bottomrule
\end{tabularx}
\end{center}
\end{table}

The model does not reach breakeven in year two, and presenting it otherwise
would be dishonest. Breakeven under these assumptions occurs in year four at
approximately 78{,}000 paying users with retention above 65\%. The sensitivity
is dominated by conversion rate: a move from 3.2\% to 4.5\% pulls breakeven
forward by roughly fourteen months, while a move to 2.0\% pushes it beyond the
model's horizon.

Consequently the highest-leverage commercial activity is not acquisition volume
but conversion, and the highest-leverage conversion mechanism is the live stock
dial of Section~\ref{sec:ui} --- the moment a user sees their own image
transform under a stock they cannot access. This is a case where a product
decision and a business decision coincide exactly, and it should be instrumented
and optimized accordingly.

\subsection{Go-To-Market}

Three sequenced motions:

\paragraph{Credibility first (pre-launch, months 19--24)} Seed the product with
the professional community: colorists, cinematographers, film photographers with
audiences. The objective is not reach but verdict. If people who know what
printer lights are say the printer lights are right, that assessment propagates
through the enthusiast tier with far more force than advertising.

\paragraph{Education as marketing (launch onward)} The explanation content built
for FR-20 and the long-press affordances is genuinely useful standalone material.
Published as short articles and videos --- what a characteristic curve is, why
pushing film does not recover shadows, what halation actually is --- it
functions simultaneously as marketing, as onboarding, and as a demonstration
that the team knows the subject. This is a durable channel that does not decay
like paid acquisition.

\paragraph{Comparison content (post-launch)} Side-by-side renders against real
film scans, published with the source files. This is only available to a product
confident in its fidelity, which is precisely why it is persuasive.

\subsection{Risks to the Commercial Thesis}

The commercial risks are distinct from the engineering risks of
Section~\ref{sec:risks} and are enumerated there alongside them. The dominant
one is stated here because it shapes strategy: \textbf{the population that can
perceive the difference may be smaller than assumed}. If most users cannot
distinguish a physically simulated render from a good LUT at phone viewing size,
the product's entire investment thesis is mispriced.

The mitigation is measurement, early. The perceptual evaluation of
Section~\ref{sec:validation} should be extended during Phase 1 to include a
non-expert cohort, testing not whether they can identify real film but whether
they \emph{prefer} the physically simulated render to a strong LUT baseline. A
preference rate near chance among non-experts would be a signal to reposition
toward the professional tier and price accordingly, and it is far better to
learn that in month seven than in month twenty-four.
